\section{INTRODUCTION}
    \subsection{Background of the project}
International School Manila (ISM) handles a wide range of academic and administrative reports that must remain authentic, traceable, and tamper-proof. These records underpin accountability, accreditation, and academic credibility. With ISMs Standard process, reports are produced by different applications and uploaded directly to on-premise servers or cloud storage through APIs that create each student records. This process lacks immutable version tracking and an audit trail. Files or metadata can be modified or replaced without a durable record of who changed what and when. Since there is no cryptographic binding between a stored file and its database record, undetected alterations are possible and unverifiable. When reports are shared or reviewed, external like colleges and universities parties have no independent way to prove that the files match the originals released by ISM. This requires multiple email exchange between ISM and the external organization.The absence of tamper-evidence and end-to-end traceability weakens ISM’s data governance and institutional trust.

The volume of existing records gives this risk concrete institutional consequences. Between 2014 and 2025, ISM’s student report repositories (covering report cards, CAS reports, Measures of Academic Progress results, and related documents) already contain \textbf{112,612 files}, with a further \textbf{16,108 files} generated by Finance. As of 2025, ISM has \textbf{2,445} active students, each of whom accumulates dozens of academic and financial records even before alumni and long term archives are counted. In this setting, even small gaps in authenticity, provenance, or version control can affect tens of thousands of files at once.

In this environment, the people inside ISM who rely on these reports for decisions and reviews cannot follow a clear, tamper-evident trail of what has happened to a file over time. Instead of checking a cryptographic record that proves the report is unchanged, they have to trust that whatever is in front of them is the “correct” version the institution intends. Over time, this kind of reliance on institutional assurances rather than mathematically verifiable proof has been shown to create both operational and compliance risk  \citep{Cflow2025,Encrypthos2025,NRC2005,SimpleButNeeded2025}. Hence, uncertainty exists during audit time and accreditation review, as well as interdepartmental review. Especially as multiple copies of the same file can exist in different systems, the process becomes slow, manual, and IT-dependent. 


Across many sectors, including universities, hospitals, and financial organizations, teams use partial remedies such as cloud storage versioning, offline checksums, and IPFS pinning services to improve recovery and traceability of files. IPFS (InterPlanetary File System) is a decentralized, peer to peer protocol that stores and shares data using content addressed identifiers, so that each file is located by a cryptographic hash rather than a server location \citep{IPFSDocs2025,CryptoSTIPFS2025}. These tools increase resilience, but they usually operate as separate utilities outside the daily workflow of academic offices. Their assurances still depend on local logs and platform trust, not on an end to end, independently verifiable trail \citep{AWSS3Versioning2025,NIST80092r1,FilebasePinning2025,PinataPinning2025,SchneiderIPFSPinning2025,BlockchainReporterIPFS2025}.


To fill this gap, \textbf{VERACISM} (Verifiable Academic Records for ISM) establishes content-addressed, decentralized persistence for report storage and verification. Each uploaded report is scanned (ClamAV and YARA), validated (MIME/size), hashed (SHA-256), and stored via Lighthouse on IPFS with Filecoin-backed deals. Each file is assigned a unique Content Identifier (CID) which will change instantly if the content changes so that tamper evidence and public verification are instant. A .NET 8 Web API supervises secure upload, and policy checks; MSSQL provides tenant aware metadata and immutable audit events, with strict RBAC/ABAC. Verification happens at a portal level with QR or CID lookups that cause the content to be re-hashed and checked against evidence of deals on Filfox, so that independent verification is possible without ISM credentials. Baseline integrity, permanence, and independent verification are completely satisfied in terms of wholely content addressed and decentralized storage, and the code of a block chain will only be used in future versions if further requirements arise (for example, on chain rules for access, payments and cross-application automation). 

VERACISM utilizes a \textbf{hybrid architecture} in which on-premise control combines with decentralized permanence: ISM takes ownership of identity, access rules, metadata, and operational audit trails within its own trusted environment, while content integrity and long-term durability are anchored to public decentralized infrastructure (IPFS/Filecoin) for independent verification. This model allows full flexibility for existing workflows and controls for compliance, but provides globally verifiable proof of integrity which does not rely on the systems of ISM to be trusted. 

VERACISM is designed to replace implicit trust with explicit proof, to support long lasting and visible record keeping, and to make audits easier to perform. By combining cryptographic integrity, decentralized data storage, and an outer layer of verification, the proposed system aims to strengthen institutional credibility, reduce the overhead of verification checks, and improve protection of reports against silent or unauthorized modification.


    \subsection{\\Objectives}
        The proposed system aims to modernize and secure the creation, storage, and independent verification of ISM academic reports by leveraging content-addressed, decentralized infrastructure (IPFS/Filecoin via Lighthouse). The primary objective is to prevent tampering and streamline verification through CID- and QR-based checks that provide cryptographic, tamper-evident proof of integrity without requiring smart contracts.
        
    \subsection{Specific Objectives}
\begin{enumerate}

  \item Implement a file security checker and file-type validation.
  \item Design and develop the web application with role-based access control:
  \item Establish a threat model and mitigations: input sanitization, rate limits, CSRF, key rotation, least privilege.
   \item Define and execute integration, security, and performance tests with clear success criteria, including correct end-to-end handling of reports across integrated systems, no critical or high-severity security findings, and response times and throughput that remain within agreed service levels under expected load.

  \item Provide runbooks for deployment, keys, backup/restore, monitoring, incidents, and compliance.
\end{enumerate}

    

    \subsection{Scope and Limitations}
        The system provides immutable storage, verification, and controlled retrieval of finalized reports with cryptographic proof of integrity and a complete audit trail. It does not generate report content or replace existing academic systems. The following items summarize the project boundaries, goals, and benefits.
        \subsubsection{Scope}
        \begin{itemize}

            \item Verification Portal
                \begin{itemize}
                    \item QR and CID portals for independent verification without ISM credentials
                \end{itemize}
                
            \item Document Security Feature
                \begin{itemize}
                    \item Pre-upload security controls: ClamAV, YARA rules, MIME and size validation, quarantine workflow
                    \item Content-addressed storage using Lighthouse with IPFS/Filecoin, with CID issuance and persistence monitoring
                    \item SHA-256 hashing and CID revalidation during verification
                \end{itemize}
                
            \item Events/Access Management
                \begin{itemize}
                    \item REST APIs for upload, metadata retrieval, verification, and controlled download
                    \item Tenant-aware metadata and immutable audit events in MSSQL
                    \item Role-based and attribute-based access control with scoped API keys
                    \item Scheduled health checks for CID availability, replica status, and retrieval latency
                    \item Structured logging, alerting, and runbooks for operations across Dev, UAT, and Production
                \end{itemize}
                
            \item Integration
                \begin{itemize}
                    \item Integration with existing ISM applications through approved service accounts
                \end{itemize}


        \end{itemize}


           
             
        \subsubsection{Limitations}
            \begin{itemize}
              \item Report authoring, editing, grading, or workflow approvals
              \item Replacement of SIS or LMS and non-records business functions
              \item Smart contracts, on-chain access rules, or payment logic
              \item Student self-service uploads
              \item General file sharing, public website hosting, or email/notification design
              \item Semantic content moderation beyond malware and known bad pattern checks
              \item User crypto wallet management or token handling
              \item The system currently lacks an automated retry mechanism for transient
                    Lighthouse storage failures; a document that passes scanning may not be
                    anchored to IPFS if the Lighthouse service is temporarily unavailable.
              \item The \texttt{VerificationLog} table is not yet written during public CID
                    verification, meaning per-check verification history is not persisted.
              \item The audit log UI does not yet enforce tenant-level scoping; audit metadata
                    visibility is currently limited to the Admin role.
              \item Performance evaluation was scoped to the public verification endpoint at
                    100 virtual users; a broader load profile covering the full scan-upload
                    pipeline under sustained concurrency remains outside the current scope.
              \item The administrator dashboard does not yet include redirect links from
                    summary cards to detail or full report views; this usability gap was
                    identified during UAT and is planned for the next development phase.
            \end{itemize}